%%%%%%%%%%%%%%%%%%%%%%%%%%%%%%%%%%%%%%%%%%%%%%%%%
\subsection{Generative Models for Image Synthesis}\label{sec:lit_generative_models}

Deep generative image models split historically into two conceptual families. \textit{Generative Adversarial Networks} (GANs) \cite{goodfellowGenerativeAdversarialNetworks2014} pit a generator against a discriminator in a minimax game: the generator learns to map noise to data samples, the discriminator learns to distinguish generated from real samples, and at the (theoretical) global optimum the generator's distribution exactly matches the data distribution and the discriminator can no longer do better than chance. GANs require no Markov chain or iterative inference at sample time -- a single forward pass through the generator suffices -- but the adversarial objective is notoriously unstable to optimize, and the original paper already flags a specific degenerate failure mode (the generator collapsing many noise values onto a single output) that later work spent years addressing architecturally. \textit{DCGAN} \cite{radfordUnsupervisedRepresentationLearning2016} showed that a specific combination of architectural constraints -- strided convolutions in place of pooling, batch normalization, and asymmetric activation choices between generator and discriminator -- stabilizes convolutional GAN training well enough to learn useful, transferable image representations. \textit{StyleGAN} \cite{karrasStyleBasedGeneratorArchitecture2019} then re-designed the generator itself around a mapping network that produces an intermediate latent code applied at every convolutional layer via adaptive instance normalization, plus explicit per-pixel noise injected at each layer to model stochastic fine detail (hair, freckles) independently of global attributes (pose, identity); this style-based design both improved measured image fidelity (Fréchet Inception Distance) substantially over the traditional feed-forward generator and made the synthesis process more controllable and interpretable through style mixing.

\textit{Diffusion models} approach image synthesis from a different direction: rather than learning a single feed-forward noise-to-image mapping, they learn to reverse a fixed Markov chain that gradually corrupts an image into noise. \textit{Denoising Diffusion Probabilistic Models} (DDPMs) \cite{hoDenoisingDiffusionProbabilistic2020} formalize this as training a neural network to predict the noise added at each step of the forward corruption process, optimizing a simplified variational bound that in practice yields substantially better sample quality than the full likelihood objective; on unconditional CIFAR-10 this already matched or exceeded the FID of contemporary GANs (FID $3.17$) without any adversarial training and its associated instability. The principal early drawback of diffusion models was sampling cost: producing one sample requires simulating the reverse Markov chain over many (e.g. $1{,}000$) sequential network evaluations. \textit{Denoising Diffusion Implicit Models} (DDIM) \cite{songDenoisingDiffusionImplicit2022} address this by generalizing the forward process to a family of non-Markovian processes that share DDPM's training objective but admit a deterministic generative process, yielding samples $10$--$100\times$ faster in wall-clock time with comparable quality, while also enabling meaningful latent-space interpolation. \textit{Classifier-free guidance} \cite{hoClassifierFreeDiffusionGuidance2022} separately addressed the quality/diversity trade-off: by jointly training a single network as both a conditional and an unconditional diffusion model (randomly dropping the conditioning signal during training) and linearly extrapolating between the two score estimates at sampling time, it reproduces the fidelity gains of classifier guidance without requiring a separate classifier trained on noisy intermediate images -- a technique now used by essentially all text-conditioned diffusion models.

Bringing diffusion models to high-resolution, text-conditioned synthesis required a further architectural shift. \textit{Latent Diffusion Models} (the basis of \textit{Stable Diffusion}) \cite{rombachHighResolutionImageSynthesis2022} run the diffusion process not in pixel space but in the compressed latent space of a separately pretrained perceptual-compression autoencoder, and introduce cross-attention layers into the denoising U-Net to admit arbitrary conditioning modalities, including text; this reduces training and inference cost by roughly an order of magnitude relative to pixel-space diffusion while retaining competitive or state-of-the-art synthesis quality. \textit{SDXL} \cite{podellSDXLImprovingLatent2023} scales this recipe up substantially -- a $2.6\,\mathrm{B}$-parameter U-Net (versus $860$--$865\,\mathrm{M}$ for Stable Diffusion 1.x/2.x), a dual text-encoder conditioning scheme combining \textit{OpenCLIP ViT-bigG} and \textit{CLIP ViT-L}, additional micro-conditioning on the original (pre-crop, pre-resize) image size and crop offsets to eliminate cropping and low-resolution artifacts, and multi-aspect-ratio training -- and was shown in human preference studies to clearly outperform its predecessors. \textit{Stable Diffusion 3} \cite{esserScalingRectifiedFlow2024} departs further still, replacing the U-Net backbone with \textit{MM-DiT}, a multimodal diffusion transformer that maintains separate weight streams for image and text tokens with bidirectional attention between them, combined with a rectified-flow formulation (straight-line noise-to-data paths) and an improved, non-uniform timestep sampling distribution; it scales predictably up to $8\,\mathrm{B}$ parameters and outperforms SDXL, SDXL-Turbo, and DALL-E 3 on prompt adherence and human preference. Finally, \textit{Latent Consistency Models} (LCM) \cite{luoLatentConsistencyModels2023} show that a pretrained latent diffusion model can be distilled into a consistency model that directly predicts the solution of the underlying probability-flow ODE, collapsing sampling from dozens of steps to as few as $2$--$4$ (or even $1$) with only a fraction of the original training compute -- illustrating that inference cost, not just sample fidelity, is an actively optimized axis in this literature.

A second axis along which these models differ, orthogonal to the GAN/diffusion architectural split, is how much natural-language information they can actually condition on. Open local diffusion checkpoints inherit their text-conditioning limits from whichever text encoder they use: \textit{Stable Diffusion} 1.x/2.x and SDXL condition primarily on \textit{CLIP} \cite{radfordLearningTransferableVisual2021}, a contrastively-trained image/text encoder pretrained on $400$ million (image, text) pairs whose transformer text branch caps its input sequence at $76$ tokens for computational efficiency -- a hard architectural ceiling on prompt length, not a soft practical guideline. Later open models incorporate an additional, much larger-capacity text encoder: SD3 and \textit{FLUX.1} \cite{labsBlackForestLabs} pair CLIP-family encoders with \textit{T5} \cite{raffelExploringLimitsTransfer2023}, an encoder-decoder Transformer originally trained on the $\sim 750\,\mathrm{GB}$ Colossal Clean Crawled Corpus (C4) with a maximum pretraining sequence length of $512$ tokens; used standalone as a conditioning branch, T5's encoder gives SD3/FLUX-family models access to prompts roughly an order of magnitude longer than CLIP's, and SD3's own ablation confirms this measurably improves performance on complex, highly detailed prompts and on rendering legible text, though it has only a small effect on general aesthetic quality. In parallel, large proprietary providers moved in the opposite direction from these openly-released, self-hostable checkpoints: text-to-image and text-to-image-editing APIs such as \textit{Gemini} image generation and the \textit{GPT-image} family (whose lineage traces to \textit{DALL-E} \cite{openaiDALLECreatingImages2021}) are offered exclusively as hosted services, with proprietary architectures and undisclosed text-encoder budgets that in practice accept prompts far longer than any locally-hosted checkpoint can consume.

This split between proprietary API-based generation and open, locally-run diffusion checkpoints, and specifically the accompanying difference in usable prompt length, is not a peripheral detail for this thesis: it became a controlled experimental factor in the model comparison underlying \Cref{sec:results_synthetic_data_practicability}. As documented in \texttt{docs/synthetic-model-comparison/05\_prompt-strategy-and-length-limits.md}, the production prompt template used with the proprietary API models runs to roughly $1{,}300$ words for training images and $2{,}700$ words for test images, while the open local checkpoints evaluated in this thesis have hard token ceilings that are one to two orders of magnitude smaller: $\mathit{SDXL}$-family models (via their CLIP text encoder) are limited to $\sim 75$ usable CLIP tokens, the \textit{SD\,3.5} family to $\sim 256$ T5 tokens, and \textit{FLUX.2-klein-9B} and \textit{Qwen-Image} (via their respective Qwen-family text encoders) to $\sim 512$ tokens. Because a prompt written for a $32\,\mathrm{k}$-token-capable API model cannot simply be reused verbatim for a $75$-token-capable local model without changing what information survives, this thesis's own synthetic-data generation pipeline treats prompt length as an explicit, separately-controlled experimental regime -- a \textit{compressed} regime that gives every model an identical, tightly-budgeted prompt to isolate model quality from prompt capacity, and a \textit{maxlen} regime that instead fills each model's own token budget with genuine per-image content -- rather than treating "the prompt" as a fixed quantity that transfers unproblematically across the API/open-checkpoint divide.

%%%%%%%%%%%%%%%%%%%%%%%%%%%%%%%%%%%%%%%%%%%%%%%%%
\subsection{Synthetic Data for Object Detection and the Sim-to-Real Gap}\label{sec:lit_sim_to_real}

Long before diffusion-based image generation existed, the computer vision and robotics communities had already established that models trained purely on synthetic imagery can transfer to real-world perception tasks, provided the synthetic data is constructed with transfer in mind. \textit{Domain randomization} \cite{tobinDomainRandomizationTransferring2017} demonstrated this in its starkest form: an object-localization network trained exclusively on low-fidelity rendered images with randomized, deliberately non-photorealistic textures, lighting, camera pose, and distractor objects achieved $1.5\,\mathrm{cm}$ real-world localization accuracy with no real-image pretraining at all. The underlying hypothesis -- that with sufficient variability in the simulated distribution, the real world simply appears to the model as one more random variation rather than a distinct domain -- explicitly runs counter to the intuition that closing the sim-to-real gap requires photorealism; the paper notes that prior attempts to transfer purely via realistic RGB rendering had in fact had only limited success. \textit{Tremblay et al.} \cite{tremblayTrainingDeepNetworks2018} extended domain randomization from simple localization to full 2D object detection, rendering non-photorealistic cars and geometric "flying distractor" objects over random real-photo backgrounds and evaluating on the real-world KITTI benchmark. Their domain-randomized detector was competitive with, and in two of three tested architectures outperformed (e.g. SSD: $46.3$ vs. $36.1$ AP@0.5), a detector trained on \textit{Virtual KITTI}, a substantially higher-fidelity synthetic dataset purpose-built to visually match the KITTI test distribution -- indicating that forcing a network to ignore rendering-level texture and lighting cues can matter more for real-world transfer than matching those cues in the first place. Fine-tuning the domain-randomized detector afterward on a modest amount of real data improved results further, beyond training on real data alone.

Generative models introduce a different, and in some respects harder, version of the sim-to-real problem, since the synthetic distribution is now defined implicitly by a learned generator rather than by an explicit, controllable rendering pipeline. \textit{He et al.} \cite{heSyntheticDataGenerative2023} conducted the first systematic study of this question for text-to-image diffusion output (using GLIDE), finding that synthetic images can measurably improve zero-shot classification when layered on top of a strong pretrained feature extractor (a $4.31$-point average top-1 gain across $17$ datasets with a frozen CLIP ResNet-50 backbone, and $2.90$ points with a CLIP ViT-B/16 backbone, using language-diversified prompts and CLIP-based filtering to control noise), but that this benefit collapses when the same synthetic data is instead used to train a classifier from scratch: a $50\mathrm{k}$-image synthetic training set underperformed real data by a wide margin, requiring roughly five times as many synthetic images to match the accuracy obtainable from real images of the same classes. This is direct evidence of a domain gap between generated and real imagery that a strong pretrained representation can partially absorb, but that raw supervised training cannot. \textit{Azizi et al.} \cite{aziziSyntheticDataDiffusion2023} showed that this gap can be substantially narrowed, but only by fine-tuning the generator itself on the target domain rather than merely prompting an off-the-shelf model: an \textit{Imagen} diffusion model fine-tuned on ImageNet reached a state-of-the-art Classification Accuracy Score of $69.24\%$ (models trained solely on its generated images, evaluated on the real ImageNet validation set), which the authors describe as greatly narrowing -- though, on the accuracy figures reported, not fully closing -- the remaining gap to real-data training. Critically, the same paper reports that naively prompting an off-the-shelf, non-fine-tuned diffusion model with short class-name prompts produces a "very poor" classifier, and cites two contemporaneous studies that found pretrained Stable Diffusion output did not improve ImageNet accuracy at all when added to real training data -- underscoring that raw visual fidelity of the generator is not sufficient on its own, and that alignment between the generator's learned distribution and the target label space and appearance distribution is the operative bottleneck. \textit{DA-Fusion} \cite{trabuccoEffectiveDataAugmentation2025} takes an intermediate position: rather than generating synthetic images from scratch, it edits real images semantically using a pretrained Stable Diffusion checkpoint via image-to-image splicing (SDEdit), with per-class textual-inversion embeddings learned from as few as one labeled example so that the generator can render concepts, including fine-grained categories, that lie outside its own training vocabulary; this real-image-anchored augmentation outperformed both standard geometric augmentation and prior diffusion-based augmentation by up to $10$ accuracy points in few-shot classification, with the largest gains on the most fine-grained concepts -- but it works precisely because it constrains generation to remain close to a real reference image rather than sampling unconstrained from the generator's prior.

Across this body of work, three recurring failure modes motivate treating synthetic imagery as a supplement whose reliability must be actively verified rather than assumed: a persistent \textit{texture and lighting domain shift} between rendered or generated pixels and real sensor imagery, even after substantial fine-tuning effort (Azizi et al.); a \textit{label- and context-alignment} gap in which a generator's default visual associations for a class or scene diverge from the target task's actual distribution unless explicitly corrected (He et al.'s from-scratch results; Azizi et al.'s off-the-shelf failure case); and, at the level of what training on generative output can compound over time rather than within a single generation, \textit{model collapse} \cite{shumailovAIModelsCollapse2024} -- the theoretically and empirically demonstrated degeneration of a generative model's learned distribution when it is recursively retrained, indiscriminately, on data sampled from itself or its predecessors, first losing distributional tails and eventually collapsing toward a low-variance, degenerate output distribution. None of these failure modes is a purely theoretical concern for this thesis, and together they are the direct motivation for treating the practice of never judging a model on synthetic images alone (\Cref{sec:evaluation_framework}) as a methodological necessity rather than a formality: a detector whose headline metric was computed on synthetic imagery alone could look artificially strong precisely because it is being evaluated on the same distribution whose domain shift from real imagery, alignment gaps, or (in a repeated-generation pipeline) collapse risk are exactly what is in question. The mixed real-and-synthetic evaluation protocol adopted in this thesis, with the real-only breakout always reported alongside it and a real-vs-synthetic delta actively monitored as a watchdog, is a direct methodological response to this literature.
